\documentclass{subfiles}
\begin{document}
    Our discussion on RB deletion will assume the following:
        for any node \(y\), \(x.r\) denotes \(x\)'s right child,
        similarly \(x.l\) denotes the left child,
        with \(x.p\) we indicate \(x\)'s parent and with \(x.c, x\) 's color.
        Additionally, we denote with \(z\) the node to remove, 
            with \(y\) the node the replaces the deleted node and we let it always point to \(z\),
            and with \(x\) the node filling \(y\)'s position after this has been moved.

    As for insertion, we mainly focus on the properties that may fail upon deletion,
        and explain how these are fixed. 

    First, if \(y\) was the root and a red child replaces it, \Cref{prop:rb-prop-2} fails.
        Second, if \(x\) and \(x.p\) are both red, then \Cref{prop:rb-prop-4} fails.
        Third, moving \(y\) may cause \Cref{prop:rb-prop-5} to fail.
        We can correct the violation of \Cref{prop:rb-prop-5} by saying that node x,
        now occupying y’s original position, has an ``extra'' black.
        That is, if we add 1 to the count of black nodes on any simple path
        that contains x, then under this interpretation, property 5 holds. When we remove
        or move the black node y, we “push” its blackness onto node x. 
        The problem is that now node x is neither red nor black, thereby violating property 1.

    \subfile{../TikZ/Procedure - RB deletion.tex}
    
    With respect to \autoref{proc:RB-deletion}, which shows the pseudo code to implement RB deletion,
        we show how \lstinline{RB-Delete-Fixup} restores each of the RB properties.
        \begin{description}
            \item [Case 1: \(\mathbf{z}\)'s sibling \(\mathbf{w}\) is red.]
                We know that \(w\)'s must be black, therefore we can exchange the colors of 
                \(w\) and \(x.p\) and apply a left rotation on \(x.p\).

            \item [Case 2: \(\mathbf{z}\)'s sibling \(\mathbf{w}\) is black, and so are its children.]
                Since both \(w\)'s children are black, and so are \(z\) and \(w\),
                we can simply pulling away a black from both, leaving \(z\) singly black and \(w\)
                red.

            \item [Case 3: \(\mathbf{z}\)'s sibling \(\mathbf{w}\) is black,
                \(\mathbf{w.l}\) is red and \(\mathbf{w.r}\) is black.]
                We can switch the colors of \(w\) and its left child \(w.l\)
                    and then perform a right rotation on \(w\).

            \item [Case 4: \(\mathbf{z}\)'s sibling \(\mathbf{w}\) is black, 
                and \(\mathbf{w.r}\) is red.]
                By making some color changes and performing a left rotation on \(x.p\),
                we can remove the extra black on \(x\), making it singly black.
        \end{description}
        Upon termination \Cref{prop:rb-prop-2} may still fail,
            which is fixed at the end of \lstinline{RB-Delete-Fixup}.
            \subfile{../TikZ/Procedure - RB delete fixup.tex}
\end{document}
